\documentclass[10pt,a4paper]{amsart}
\usepackage[margin=19mm]{geometry}
\usepackage{amsmath,amssymb,mathtools}
\usepackage[scr=rsfs,cal=euler]{mathalfa}
\usepackage{xfrac}
\usepackage[unicode,hidelinks,bookmarks=false]{hyperref}
\newcommand{\ie}{i.e.\ }
\newcommand{\cf}{cf.\ }
\newcommand{\resp}{respectively\ }
\newcommand{\Subsection}{\subsection}
\providecommand{\nobreakdash}{\nobreak\hbox{-}}
\setlength{\emergencystretch}{2em}
\multlinegap0pt
\usepackage{bm}
\usepackage{bbm}
\usepackage{dsfont}
\usepackage{xspace}
\usepackage{cancel}
\usepackage{mathtools}
\usepackage{amsthm}
\makeatletter
\@ifundefined{theorem}{\newtheorem{theorem}{Theorem}}{}
\@ifundefined{lemma}{\newtheorem{lemma}[theorem]{Lemma}}{}
\makeatother
\newcommand{\C}{\mathbb{C}}
\newcommand{\Ga}{\Gamma }
\newcommand{\Om}{\Omega }
\newcommand{\Z}{\mathbb{Z}}
\newcommand{\ep}{\epsilon }
\title[The scaling limit of fair Peano paths]{The scaling limit of fair Peano paths}
\author{Nathan Albin, Joan Lind,, Pietro Poggi-Corradini}
\date{Source: arXiv:2106.10376v3 (CC BY 4.0)}
\begin{document}
\maketitle

\noindent\emph{Source and licence.} The scaling limit of fair Peano paths. Version arXiv:2106.10376v3 (CC BY 4.0), distributed under the Creative Commons Attribution 4.0 International licence, as recorded in the arXiv licence metadata for this exact version and shown on the abstract page https://arxiv.org/abs/2106.10376v3. Full rights notice: licenses/arxiv-2106.10376-CC-BY-4.0.txt.\par\smallskip

\noindent\textit{Editorial scope.} Counted unit: the Lemma whose source label is \texttt{lem:planar-grids} in the article (source0.tex lines 1108--1110, source label \texttt{lem:planar-grids}), delivered together with the article's own complete original proof of it (source0.tex lines 1113--1178, one \texttt{proof} environment of 66 lines). No mathematics is written, repaired or re-derived: statement and proof are the article's own text line for line. The proof is detached from the statement in the source: the article states the result at lines 1108--1110 and prints its proof at lines 1113--1178, and the proof environment's own header names the statement, so the association is the article's own. Required context retained uncounted: none. Every definition, display and label of those lines is retained unchanged. Omitted from this package and not redistributed: 1--1107, 1111--1112, 1179--1229. Nothing inside a retained excerpt is omitted or altered: every retained body is delivered exactly as the article's lines are written, so no omission receipt of any kind accompanies this package. Citations: the retained excerpts carry no citation command at all, so no bibliography entry of the article is transcribed and no reference is redistributed. Reference adaptation: none was needed and none was made; every reference inside the delivered unit resolves inside it, so no unresolved reference or citation is produced. Printed slips kept literal: no printed word, symbol, hypothesis or number of the counted statement or of its proof was altered. Layout and class: the article's own class is \texttt{elsarticle} with packages including enumerate, amsmath, amsthm, amsfonts, geometry, graphicx, amssymb, float, hyperref, bbm, amscd, algorithm, algorithmic, latexsym; the wrapper is the lane's pinned \texttt{amsart} editorial wrapper at 10pt on A4, which restates the theorem-like environments the retained text needs and adds the article's own macro definitions that the retained text needs (\texttt{\textbackslash C}, \texttt{\textbackslash Ga}, \texttt{\textbackslash Om}, \texttt{\textbackslash Z}, \texttt{\textbackslash ep}), with extra wrapper packages (bm, bbm, dsfont, xspace, cancel, mathtools). The delivered numbering is the unit's own. Assets: no retained span contains a figure environment, a graphics-inclusion command, a PDF-page inclusion, a table or a TikZ picture; no image file is copied into this package, the package declares no asset dependency and no image file is redistributed with it. Licence: version arXiv:2106.10376v3 of this article is distributed under the Creative Commons Attribution 4.0 International licence, as recorded in the arXiv licence metadata for this exact version and shown on the abstract page https://arxiv.org/abs/2106.10376v3. A full rights notice is delivered at licenses/arxiv-2106.10376-CC-BY-4.0.txt. Characters: every retained excerpt is pure ASCII, so the delivered engine, XeLaTeX with its default Latin Modern text font, reports no missing character.
\begin{lemma}\label{lem:planar-grids}
 The $m$-by-$n$ planar grid $G_{m,n}$ is strictly dense.
\end{lemma}

\begin{proof}[Proof of Lemma \ref{lem:planar-grids}]
First note that:
\[
|V(G_{m,n})|=mn\qquad\text{and}\qquad |E(G_{m,n})|=(m-1)n+m(n-1)=2mn-m-n
\]
So
\[
\theta(G_{m,n})=\frac{2mn-m-n}{mn-1}=2-\frac{m+n+2}{mn-1}
\]
In particular,
\[
\frac{\partial \theta(G_{m,n})}{\partial m}=-\frac{(mn-1)-n(m+n+2)}{(mn-1)^2}=\frac{(n+1)^2}{(mn-1)^2}>0,
\]
and by symmetry the same holds for the derivative in $n$, meaning that the $1$-density of rectangular grids is strictly increasing with the size of the grid, is always less than $2$, and in fact tends to $2$ as $m,n\rightarrow\infty$.

Now let $H$ be a proper connected vertex-induced subgraph of $G=G_{m,n}$. We first consider the smallest rectangular grid containing $H$. Namely, let
\[
i_{min}:=\min\{i:\ \exists j, (i,j)\in V(H)\}
\]
and similarly define $i_{max}, j_{min}$, and $j_{max}$. Then, define the translated grid
\[
\tilde{G}:=(i_{min},j_{min})+G_{i_{max}-i_{min},j_{max}-j_{min}}.
\]
If $H=\tilde{G}$, then $\theta(H)=\theta(\tilde{G})< \theta(G)$ and we are done. If $H$ is a proper subgraph of $\tilde{G}$, then since $\theta(\tilde{G})\le \theta(G)$, it will be enough to show that $\theta(H)<\theta(\tilde{G})$. In particular, without loss of generality we can assume that $\tilde{G}=G_{m,n}$.


As an embedded planar graph $G_{m,n}$ has $(m-1)(n-1)$ bounded faces. Whenever such a face has all four vertices in $V(H)$ we fill it in so that $H$ becomes a connected compact set $\hat{H}$ in the plane. In particular, the complement of $\hat{H}$ is an open set consisting of finitely many bounded components $\Om_j$, $j=1,\dots,k$, and one unbounded component $\Om_0$. The boundary of each one of these components can be parametrized by a curve $\Ga_j$, $j=0,\dots,k$. This can be seen by thickening $\hat{H}$ a little bit, for instance by defining
\[
\hat{H}_\ep:=\bigcup_{z\in\hat{H}}\{w\in\C: |w-z|\le \ep\}.
\]
For each $\ep$, the boundary of $\hat{H}_\ep$ consists of $C^1$-smooth Jordan curves $\Ga_{j,\ep}$, $j=0,\dots,k$. If we parametrize each $\Ga_{j,\ep}$ with its arc-length parametrization, then, as $\ep$ tends to $0$, they converge uniformly to the curves $\Ga_j$.
 
We begin by looking at the boundary of the unbounded component $\Om_0$ and think of it as being parametrized counter-clockwise. In particular, $\Ga_0$ can be taken to be piecewise linear so that the derivative $\Ga_0^\prime$ is well defined away from the nodes of $\Z^2$.
Since $\Ga_{0,\ep}$ bounds the simply connected domain $\Om_{0,\ep}\cup \{\infty\}$, by the argument principle, the change in argument for the derivative $\Ga_{0,\ep}^\prime$, as the curve completes one full loop, is $2\pi$. As $\ep$ tends to $0$, this property is inherited by $\Ga_0$, as long as the argument of the derivative $\Ga_0^\prime$ is properly defined.
At any moment when $\Ga_0$ is not at a node of $\Z^2$, the right hand of a walker traveling along with $\Ga_0$ will always be touching $\Om_0$. In particular, the only changes in the argument of $\Ga_0^\prime$ that are allowed when $\Ga_0$ passes through a node of $\Z^2$ are 
\[
0,\frac{\pi}{2},\pi,\text{ and, }-\frac{\pi}{2}.
\]
In words, either the walker goes straight, turns $90^\circ$ left, does a $180^\circ$ turn in the positive direction, or turns $90^\circ$ right. This can be verified by looking at the unbounded component $\Om_{0,\ep}$ of the thickened $\hat{H}_\ep$. Indeed, if $\Ga_0$ were to do a $180^\circ$ turn in the negative direction, then while walking along $\Ga_{0,\ep}$ the right hand would be touching the thickened neighborhood of an edge, and this neighborhood would disappear as $\ep$ tends to $0$, so that edge would not be part of the boundary of $\Om_0$, which leads to a contradiction. 

Now assume $v\in \Z^2$ is a node where $\Ga_0$ makes a right turn ($90^\circ$  in the negative direction), and let $u\in\Z^2$ be the node visited by $\Ga_0$ just before $v$, and $w\in\Z^2$ the one visited just after $v$.
Then, $u,v,w$ bound a face $f$ of $G_{m,n}$ and the fourth corner $t$ does not belong to $H$. To see why, assume by contradiction that all four corners are in $V(H)$, then the face would be contained in $\hat{H}$. But then the right hand would not be touching the unbounded component $\Om_0$ in this case.
Also, note that since $H$ is vertex-induced the two sides of $f$ incident at $t$ also do not belong to $E(H)$. 

Now, if we add $t$ to $V(H)$, to get a new vertex-induced graph $H^\prime$. Since one of the coordinates of $t$ is equal to one of the coordinates of either $u$ or $w$, we see that $H^\prime$ is still a subgraph of $G_{m,n}$. Moreover, we are guaranteed that 
\[
E(H^\prime)\ge E(H)+2.
\] 
In particular,
\[
\theta(H^\prime)= \frac{E(H^\prime)}{V(H^\prime)-1}\ge\frac{E(H)+2}{V(H)}
\]
Since $H$ is a subgraph of the grid $G_{m,n}$, we always have $\deg_H(x)\le 4$. Moreover, at least one vertex of $H$ has $\deg_H(x)<4$. Therefore, by the Handshake Lemma,
\[
2E(H)=\sum_{x\in V(H)}\deg_H(x)<4 V(H).
\]
This implies that $\theta(H')>\theta(H)$ because
\begin{align*}
 \frac{E(H)+2}{V(H)} >\frac{E(H)}{V(H)-1}
& \Longleftrightarrow (E(H)+2)(V(H)-1)> E(H)V(H)\\
& \Longleftrightarrow 2V(H)> E(H).
\end{align*}
If $H^\prime=G_{m,n}$, we are done. Otherwise, we can repeat the same argument with $H$ replaced by $H^\prime$. This process has to end, so without loss of generality, we can assume that $\Ga_0$ never makes any right turns. In particular, the argument of $\Ga_0^\prime$ can only change by $0$, $\pi/2$, or $\pi$. So since there has to be an even number of argument changes that are positive, and the sum must be $2\pi$, we either get two changes by $\pi$ or four changes by $\pi/2$. In either case, we see that $\Ga_0$ describes the boundary of  a rectangular grid, and by assumption this grid has to be $G_{m,n}$.

Finally, consider a bounded component $\Om_j$ with $j\ge 1$. Assume that the boundary of $\Om_j$ is described in the clockwise direction by a curve $\Ga_j$, constructed as $\Ga_0$ above, using the thickened set $\hat{H}_\ep$. Once again the allowed argument changes at a grid node are $0,\pi/2,\pi$, and $-\pi/2$. Also, when walking along $\Ga_j$ the right-hand touches $\Om_j$ and the total argument change must equal $-2\pi$. In particular, there must always be at least four right turns with argument change $-\pi/2$. By repeating the argument above we see that each right turn identifies a vertex in $\Om_j$ that can be added to $H$ in a way to increase the $1$-density, and therefore after finitely many steps we get that the component $\Om_j$ has been filled in. Hence, after finitely many steps, $H=G_{m,n}$. This proves the lemma.
\end{proof}

\end{document}
